\section{Experiments}
\label{sec:experiments}

We first establish measurement fidelity, then compare direct Codex with published specialized
systems and with their shared-backbone reproductions. The trace analysis follows the three
findings in the abstract: whether additional workflow stages improve accuracy, how decomposition
preserves or loses information, and when domain conventions, task conditions, and targeted
verification help specialized systems. Figure~\ref{fig:main-results}
summarizes \StudyBenchmarkCount{} benchmark records after checking the current HF branches
on 2026-09-18. Published-reference eligibility, scored Codex variants, and paired-comparator
completion are recorded separately. Ongoing comparator
runs evaluate fixed 25\% subsets; their unfinished results remain \todo{pending}.
The Codex configuration and per-benchmark operating conditions are given in
Section~\ref{sec:setup} and Appendix~\ref{app:codex-setup}.

\subsection{Experimental rigor}
\label{sec:rigor}

\paragraph{RQ1: How accurate is the benchmark SOTA finder?}
The pipeline has audited 554 candidate benchmark records. Its historical screening artifacts
contain an 83-entry evidence shortlist and a 30-task feasibility-review inventory; the current
study includes \StudyBenchmarkCount{} benchmarks, with SSTQA replacing LongBench v2. These checkpoints
reflect continuing selection by scope, comparator evidence, testability, and evaluation conditions
(Table~\ref{tab:funnel}). They are coverage counts, not finder-accuracy estimates or counts of
benchmarks cleared for execution. To measure finder accuracy, we require an independently
adjudicated set of benchmark--system--score bindings, recording the metric, split, backbone,
method variant, and source cell for each. Exact agreement on the full binding and correctness
of each field will be reported separately, with unresolved and unavailable sources counted
explicitly. \todo{Adjudicate a fixed sample; report sample size, binding accuracy, field-level
accuracy, and uncertainty.} A high rate of evidence-backed records alone cannot establish that
all stronger systems have been found.

\paragraph{RQ2: How faithfully do the executions reproduce the intended evaluation?}
\label{sec:audit}
The 2026-09-11 alignment audit examined 25 bindings selected because they had not previously
been verified. An auditor and an adversarial verifier reviewed each binding: 308 findings were
raised, five refuted or dropped, and 303 retained, including 62 blocking, 145 material, and
96 minor findings. Twenty-two of the 25 bindings had at least one mismatch. This 88\% rate
characterizes the selected audit set; it does not estimate the fraction of published systems
that are irreproducible. Protocol, metric, and stale-record findings account for 50, 38, and
35 issues, respectively (Appendix~\ref{app:audit-breakdown}). These findings motivated repairs
and explicit declarations rather than post hoc score corrections.

Execution fidelity and reproduction of a paper's numerical result are separate measurements.
For example, the completed AgentSPEX medium run was graded by the vendored official parser,
which agreed with AgentSPEX's own verdict on all 403 questions. This verifies scorer agreement
on that run, while changing its backbone still prevents interpreting the score difference as
paper-reproduction error. \todo{Complete the per-system fidelity matrix: released components,
patches, environment, judge, split, and agreement with the reference evaluator; report numerical
reproduction error only where the paper's operating conditions are actually reproduced.}

Evaluation-contract failures are reported as reproducibility findings. In the completed
Test of Time format ablation, direct Codex scored 248/280 (88.57) without an extra answer-format
sentence and 274/280 (97.86) with it, a paired gain of 9.29 percentage points (recorded exact
McNemar $p=6.2\times10^{-6}$). A subsequent source check established that the sentence came
from the TISER baseline rather than the selected REMem-I method. The added sentence therefore
remains an ablation; 97.86 is not substituted for the main condition. The matched REMem-I evaluation on the original 280 questions is \todo{pending}. This case illustrates why a measured improvement must be attached
to the protocol that produced it, rather than presented as a general advance in agent reasoning.

\subsection{Main results}
\label{sec:main-results}

\paragraph{RQ3: Where does direct Codex stand relative to published systems?}
Figure~\ref{fig:main-results} and Table~\ref{tab:all-results} cover \StudyBenchmarkCount{}
benchmarks after incorporating the reruns. We retain \RetainedCodexCount{} scored Codex results:
\PrimaryCodexCount{} primary or explicitly scoped settings and \VariantCodexCount{} labelled
protocol variants. MedAgentsBench is a valid standalone result without a matching nine-subset
published aggregate. A missing comparator run does not invalidate a Codex measurement.
RareBench's failed-network run remains diagnostic, with its primary score pending.
The LongBench v2 measurement is retained separately in Appendix~\ref{app:supplementary-results};
it is outside the current suite and its aggregate counts.

\PaperComparisonCount{} settings support descriptive published-score comparisons:
\PaperHigherCount{} higher point estimates (\PaperHigherPercent\%) and
\PaperLowerCount{} lower (\PaperLowerPercent\%), with no exact ties. This includes the completed
150-question ChemBench open-web pilot, explicitly compared with a larger published evaluation;
excluding that pilot yields seven higher and three lower estimates across 10 settings.
These are not statistical win/loss/equivalence rates. Scored variants stay visible but do not
enter this tally. Published backbones and settings still differ, so the gaps do not isolate
scaffold effects. OOLONG's refreshed score is retained, but its paper reference uses a different
non-numeric scorer; it stays outside the tally pending aligned regrading or a numeric-only join.
Across the full \StudyBenchmarkCount{}-benchmark suite, the \PaperHigherCount{} higher and
\PaperLowerCount{} lower estimates account for \SuiteHigherPercent\% and
\SuiteLowerPercent\%, respectively; the remaining \PaperUnclassifiedCount{} benchmarks
(\SuiteUnclassifiedPercent\%) are unclassified because
their comparison requires an aligned protocol, a matching reference, or a valid primary run.
Unclassified results are not evidence of parity.
\todo{Refresh after the remaining primary runs; define a practically meaningful
margin for final above/below/on-par reporting. Non-significance alone is not equivalence.}

The update replaces OOLONG's incompatible staged-file result with its 192-question labels-off,
in-prompt result (73.44), and incorporates the STaRK and ELAIPBench recovery runs by task ID.
\todo{Verify the required gold-source trajectory scan for the open-web ChemBench pilot.}
Aggregation is benchmark-specific: the Biomni rerun scores 59.76 as a ten-task macro, not the
job-level micro score of 61.89; MedAgentsBench scores 50.51 across all 862 native grades.
The new SSTQA result is 77.23 on 764 English-split questions, retained as a judge and language
variant. Its primary evaluation uses the Chinese split and remains \todo{pending}; regrading
the English predictions alone would not establish that primary result. None of these additions
requires a completed specialized-system run.

The preliminary pattern varies across tasks. Codex scores 73.97 versus ARK's published 48.20
on STaRK-PRIME, and 61.09 versus RARG's 51.75 on the evaluated BRIGHT domains. Financial reasoning
has much smaller external gaps: 79.00 versus 76.00 on FinanceMath and 71.26 versus 71.17 on
MultiHiertt. Codex is below the selected published point estimates on EarthBench (59.11 versus
65.99; 65.59 from released outputs), Test of Time (88.57 versus 93.10), and ELAIPBench (37.47 versus 43.70). Tool asymmetries,
subset coverage, and missing trials remain attached to their rows, so a lower point estimate
alone does not establish a capability deficit. EarthBench's answer-key revisions and the
resulting sensitivity analysis remain separate from this historical-score comparison
(Appendix~\ref{app:result-details}).

\begin{table}[t]
\caption{Task taxonomy of the \StudyBenchmarkCount{} reported benchmarks. These mutually exclusive
presentation groups organize domain analysis; they do not replace the collection's overlapping
33-label taxonomy. Counts include scored variants and the pending RareBench rerun. No average pools these
heterogeneous metrics.}
\label{tab:domain-taxonomy}
\centering
\small
\setlength{\tabcolsep}{4pt}
\begin{tabular}{@{}p{3.1cm}rp{8.7cm}@{}}
\toprule
Task family & $n$ & Benchmarks \\
\midrule
Retrieval and long context & 5 & STaRK-PRIME, BRIGHT, SSTQA, CRAG, OOLONG \\
Scientific reasoning & 7 & SciTab, LLM-SRBench, EarthBench, MatSciBench, ChemBench,
FrontierScience Research, ELAIPBench \\
Clinical and biomedical & 3 & Biomni Eval1, MedAgentsBench, RareBench \\
Quantitative and data tasks & 4 & IndustryOR, FinanceMath, MultiHiertt, InsightBench \\
Temporal reasoning & 1 & Test of Time \\
\bottomrule
\end{tabular}
\end{table}

Table~\ref{tab:domain-taxonomy} defines the grouping for domain analysis. We additionally tag
whether success requires corpus retrieval, executable mathematical models, domain APIs,
multimodal perception, or long-context aggregation. \todo{Complete these capability tags and
report within-family paired effects and failure rates once comparator samples finish.}
The incomplete clinical and scientific rows preclude a current ranking of domain difficulty.

\paragraph{RQ4: How do the differences change on a shared backbone?}
\label{sec:likeforlike}
Table~\ref{tab:l4l} joins three completed comparator evaluations to Codex by instance ID.
FinanceMath and IndustryOR use all 200 and 100 questions. ELAIPBench now uses all 403:
the mentor re-executed the four original pre-agent failures, obtaining four valid zero-credit
outcomes, so Codex scores 151/403 and AgentSPEX-medium 184/403. The recovery records replace
the failed trials by ID; they are not appended as additional questions.

\begin{table}[t]
\caption{Paired results on the shared \texttt{gpt-5.6-sol} backbone. Scores are percentages;
$\Delta$ is Codex minus comparator in percentage points. $C$-only/$S$-only count questions
solved exclusively by Codex/the specialized system. Intervals are 95\% paired bootstrap
intervals, resampling papers for ELAIPBench and instances otherwise. $p_{\mathrm H}$ is the
item-level exact two-sided McNemar $p$-value with Holm correction across these three comparisons.
MoCA$^*$ denotes the released ablation without the unavailable baseline proposer, selector
committee, and conflict arbiter.}
\label{tab:l4l}
\centering
\scriptsize
\setlength{\tabcolsep}{3.2pt}
\begin{tabular}{@{}llrrrrlrr@{}}
\toprule
Benchmark & Comparator & $n$ & Codex & System & $\Delta$ & 95\% interval & $C$/$S$-only & $p_{\mathrm H}$ \\
\midrule
FinanceMath & MoCA$^*$ & 200 & 79.00 & 79.00 & $0.00$ & $[-4.50,4.50]$ & 10/10 & 1.000 \\
ELAIPBench & AgentSPEX-medium & 403 & 37.47 & 45.66 & $-8.19$ & $[-12.87,-3.42]$ & 34/67 & 0.0040 \\
IndustryOR & ReLoop & 100 & 73.00 & 70.00 & $+3.00$ & $[-3.00,9.00]$ & 6/3 & 1.000 \\
\bottomrule
\end{tabular}
\end{table}

On FinanceMath the published-score gain of 3.00 points becomes a tie in aggregate accuracy;
on IndustryOR, the gain falls from 5.00 to 3.00 points. Neither paired interval establishes
an advantage or equivalence. AgentSPEX retains an 8.19-point advantage at explicit medium effort.
Its separate high-effort run scores 190/403, versus 184/403 at medium; 21 questions favor high
and 15 favor medium (exact McNemar $p=0.405$). Thus the observed ELAIPBench gap persists when
reasoning effort is explicitly matched; these measurements do not establish that effort has
no effect. The FinanceMath run omitted the effort argument and is recorded under the project's
endpoint-default-medium assumption. Both FinanceMath arms use its official numeric matcher,
including its sign/scale tolerance, and both IndustryOR arms use the original answer key.

We recomputed the paired statistics from saved per-instance verdicts. The three tests are
exploratory, and benchmarks were not sampled randomly. The item-level McNemar tests do not account for within-paper dependence; Holm correction
addresses the three comparisons, not that dependence. ELAIPBench has 137 source papers;
its paper-cluster bootstrap interval remains below zero, while a cluster sign-flip check gives
Monte Carlo $p=0.0015$. The export records all 703 paired outcomes, 20{,}000 bootstrap replicates,
50{,}000 sign-flip draws, and seed 20260918. These intervals characterize variation across the
observed tasks or papers, not repeated stochastic runs of each agent.

\paragraph{Ongoing comparator evaluation.}
Five full comparator executions are recorded (FinanceMath, IndustryOR, ELAIPBench, CRAG, and
SSTQA); completion alone does not make every run eligible for a paired headline. For the remaining
rows, the 25\% manifests fix the seed and task identifiers before execution, with stratification
where declared. ChemBench instead uses a separately specified 280-item subset; Test of Time
uses the existing balanced 280-question Codex set (10\% of 2{,}800), superseding its earlier
700-question manifest. Only the ordering of that fixed set is shuffled, with seed 20260918.
Codex is scored
on the identical selected identifiers; the full-set published reference remains a separate
comparison. Completed pilot stages are useful for fidelity and cost checks, but are not final
25\% scores. \todo{Complete the sample runs and required judging; add per-row paired effects,
uncertainty, sample coverage, tokens, cost, and wall-clock time.}

\subsection{Analysis}
\label{sec:analysis}

We use the paired outcomes to locate disagreements and inspect saved intermediate artifacts,
tool results, and submitted answers. The three completed comparisons contain 130 discordant
questions: 101 in ELAIPBench, 20 in FinanceMath, and nine in IndustryOR. Below we inspect
selected discordant pairs and a successful repair on a question both agents solve; these cases
do not estimate how frequently each mechanism occurs. A trace can identify where two executions diverged; isolating the contribution of a
workflow component additionally requires a matched intervention.

\paragraph{RQ5: Do additional workflow stages improve accuracy?}
MoCA$^*$ introduces a claim catalog, role-specific claim review, market clearing, and program
synthesis, yet its FinanceMath accuracy equals direct Codex's 158/200. The tie masks ten unique
successes per arm. ReLoop's extraction, code generation, and repair workflow scores 70/100 on
IndustryOR versus Codex's 73/100, with six Codex-only and three comparator-only successes.
Together these results show that a specialized workflow need not improve aggregate accuracy
on the tested tasks; their intervals do not establish equivalence or a general penalty for
complexity. AgentSPEX's advantage on ELAIPBench also rules out interpreting these examples as
evidence that specialization is uniformly unhelpful.

ReLoop's saved stage outputs yield 69, 70, 70, and 70 correct answers after initial generation,
L1 execution verification and regeneration, L2 behavioral checks, and final repair, respectively.
The one recovered case, problem 16, initially
fails on a brittle string check for purchase timing; the repair allows the program to enforce
the extracted purchases-before-sales condition. These are successive outputs of one execution,
not controlled ablations. Later stages also leave unresolved errors: on IndustryOR
problems 52, 80, 82, and 96, ReLoop makes three repair attempts but returns no numeric objective;
Codex solves all four. On FinanceMath \texttt{validation-147}, MoCA$^*$'s initial catalog already
contains the correct three-year return and the corresponding Calmar ratio. Subsequent review
and synthesis nevertheless produce a different, incorrect calculation. These observations
motivate measuring what each stage changes, rather than using the number of stages as a proxy
for capability. \todo{Report paired model calls, executed tool calls, input/output tokens, and
latency; compare full workflows with component ablations under matched budgets.}

\paragraph{RQ6: How can decomposition lose information that Codex retains?}
ReLoop's code generator receives a summary of extracted data that omits nested field names.
For IndustryOR problem 82, the extractor records shop-count bounds under
\texttt{minimum\_number\_of\_shops}, but the generator receives only
\texttt{store\_types: list[5] of dict}. Its generated program guesses other key names and
raises a \texttt{KeyError}. The original problem text remains available; the lost information
is the schema of the extractor's own representation. Related handoff failures occur for the
staffing requirements in problem 52 and cutting demands in problem 96. Problem 52 also encounters
a missing SciPy dependency, so information loss is not its only failure source.

Information can also be present but cease to constrain the downstream solution. In FinanceMath
\texttt{validation-147}, the claim catalog identifies the three-year return as 6.2\% and computes
$6.2/10.2\approx0.6078$. The review stage marks the three-year-return claim uncertain, and the
final program instead averages returns across four distinct horizons, yielding 0.5049. Codex
uses the three-year table entry directly and returns 0.607843, which passes the evaluator.
This case exposes a failure to preserve the role of a relevant fact across stages, rather than
a failure to extract the fact initially. It illustrates how retaining the relation between
the question, table headers, and computation can matter more than adding review stages.
\todo{Test full-schema handoffs and preservation of source-linked task conditions without
changing the backbone, dependencies, or resource budget; measure repair success on all eligible
cases, including cases where the original specialized workflow succeeds.}

\paragraph{RQ7: When does domain-specific scaffolding add value?}
The comparator-only successes show concrete conditions that a useful scaffold can preserve.
On FinanceMath \texttt{validation-92}, MoCA$^*$ explicitly records that three months have elapsed
on a one-year swap and values the remaining nine months, returning 100.753663. Codex uses four
remaining quarters and returns 100.5241705. On \texttt{validation-121}, MoCA$^*$ represents the
currency units explicitly and multiplies HKD/CNY by CNY/AUD, returning 5.96844105 HKD/AUD;
Codex returns the inverse. Both specialized outputs pass the official evaluator. These examples
connect the observed advantage to retaining a time condition and resolving a quote convention,
although removing those explicit representations has not yet been tested.

IndustryOR problem 6 supplies a complementary model-convention example. Codex first solves an
integer model with objective 10755, then changes production, inventory, and backlog variables
to continuous variables and submits 10750. ReLoop retains the integer formulation and obtains
the reference objective 10755. This is a divergence in the selected model domain, not an inability
to execute the correct integer program. A targeted verification step could check variable domains
against the intended task before accepting the final model. \todo{Test convention checks,
explicit preservation of task conditions, and targeted verification separately; do not infer
their causal benefit from selected successful traces alone.}

ELAIPBench provides the strongest current aggregate advantage for a specialized workflow:
AgentSPEX has 67 unique successes versus Codex's 34. The deficit occurs in both question types:
Codex solves 35/88 single-answer questions versus 44/88 and 116/315 multiple-answer questions
versus 140/315. The stricter multiple-answer setting alone therefore cannot explain the gap.
\todo{Annotate paper evidence selection, preservation of qualifications, and verification of the
final choice across the discordant pairs; test the resulting mechanisms with controlled changes.}
The scientific pilots add a distinct measurement of tool use: S1-NexusAgent invokes no web search
on ten MatSciBench questions despite access, while its four-question Biomni pilot reaches
Ensembl, NCBI, EBI, and OpenTargets. Tool availability alone does not establish either effective
use or an accuracy benefit. The ongoing paired samples will test whether the selected tool
supplies evidence needed for the specific task.

\paragraph{Distinguishing reasoning from answer delivery.}
IndustryOR problem 28 shows why final accuracy can hide a successful computation. Codex executes
the correct objective 580000 twice, but its submitted objective field begins with the symbolic
expression $V=1.2x_3+\cdots$; the grader reads the first number, 1.2, and marks it incorrect.
This case is a delivery failure after successful computation. It remains an incorrect outcome
under the declared main protocol, and a corrected submission would be a separate intervention.
It does not establish that delivery failures are more common than reasoning errors across Codex
runs. \todo{Complete a common annotation of both arms covering evidence selection, intermediate
reasoning, execution, and final delivery; report counts and denominators before making prevalence
claims.} Protocol repairs and the Test of Time format ablation remain reproducibility evidence
(Section~\ref{sec:rigor}); they are not attributed to improved domain reasoning.
